Our experiments use public memory benchmarks (LongMemEval, LoCoMo,
and LongMemEval-V2) under their original licences; no human subjects
are involved and no new data are collected. \textsc{Structure} itself is
released as an open-source platform under a permissive licence,
including configuration files for the evaluation harness.

The event log records every context mutation, which enables auditing
and debugging but also makes the system a potentially rich source of
behavioural data. Deployments should therefore treat the event stream
as sensitive: apply retention limits proportional to operational need,
scope access to authorised users, and document retention policies to
end users of the agent. The three visibility scopes
(\textsc{user}/\textsc{workspace}/\textsc{global}) are mechanisms for
privacy control, not guarantees; misconfiguration could place
private state under a shared scope.

The system does not produce or distribute content beyond what the
underlying language model generates; the usual safety considerations
for LLM agents (misinformation, tool misuse, prompt injection through
retrieved content) apply. We recommend deployments combine
\textsc{Structure} with per-tool authorisation and content filters,
and caution that retrieved documents are themselves a prompt-injection
vector that a structured context layer does not mitigate.
